\documentclass[a4paper,11pt,fleqn]{article}%fleqn pour aligner les equations à gauche
\usepackage{/home/rweye/Documents/Overleaf/Styles/Exercice}


\newcommand{\chnum}{Chapitre 11}
\newcommand{\chtitle}{Lunette astronomique}
\newcommand{\dtype}{Exercices}
  
\begin{document}
	\maketitle
	
\begin{multicols}{2}
\section{Schématiser une lunette afocale}
La lunette astronomique de l'observatoire de Nice dispose d'un objectif de \SI{17,89}{m} de distance focale. La distance focale de l'oculaire permettant le grossissement maximal est égale à \SI{8}{mm}.
\begin{enumerate}
    \item Réaliser un schéma de la lunette astronomique afocale.
    \item Estimer la longueur de la lunette astronomique de l'observatoire de Nice.
\end{enumerate}

\section{Calculer un diamètre apparent}
Notre étoile, le Soleil, est la seule étoile qui n'apparaît pas ponctuelle depuis la Terre.\\
Calculer, en radian, l'angle $\theta$ sous lequel cet astre est vu à l'\oe il nu.\\
\textbf{Données : }\\
\textbf{Diamètre du Soleil} : \SI{1392684}{km}\\
\textbf{Distance Terre Soleil} : \SI{149 597 871}{km}

\section{Définir un grossissement}
Un astronome observe un astre à travers une lunette astronomique afocale sous un angle $\theta ' = \SI{0,20}{\radian}$. A l'\oe il nu, cet objet céleste est vu sous un angle $\theta = \SI{2E-3}{\radian}$.
\begin{enumerate}
    \item Représenter, sans soucis d'échelle, le faisceau issu d'un point objet de l'astre observé et traversant la lunette afocale. Faire apparaître les angles $\theta$ et $\theta '$ sur le schéma.
    \item Exprimer puis calculer le grossissement $G$ de cette lunette astronomique afocale.
\end{enumerate}

\section{Pouvoir séparateur de l'\oe il}
Si l'angle sous lequel l'\oe il observe deux points proches est inférieur au pouvoir séparateur de l'\oe il $\epsilon = \SI{3E-4}{\radian }$, alors les deux points sont confondus.\\
Jupiter, la plus grosse planète du système solaire, apparaît ponctuelle à l'\oe il nu depuis la Terre.
\begin{enumerate}
    \item Calculer le diamètre apparent $\theta$ sous lequel la planète Jupiter est observée depuis la Terre lorsqu'elle se situe à $D = \num{590}$ millions de km.
    \item Expliquer pourquoi cette planète apparaît ponctuelle depuis la Terre
\end{enumerate}
\textbf{Donnée : Rayon de Jupiter :} $R_J = \SI{71492}{km}$

\section{Étoiles doubles}
Deux étoiles sont vues séparées par l'\oe il si la distance angulaire qui les sépare est supérieure à \SI{3,0E-4}{\radian}. Alcor et Mizar, par exemple, constituent un couple d'étoiles visibles à l'\oe il nu dans la constellation de la Grande Ourse.\\
Mais Mizar est elle-même une étoile double invisible à l'\oe il bu car les deux composantes sont trop proches l'une de l'autre. Elles ne sont distantes que de \ang{;;14.5}.
\begin{enumerate}
    \item Exprimer en radian la distance angulaire entre les deux étoiles.
    \item En déduire le grossissement minimal de la lunette astronomique afocale permettant d'observer cette étoile double.
    \item Indiquer si cette étoile double est observable en construisant une lunette astronomique afocale à l'aide de deux lentilles minces convergentes de distances focales de \SI{20.0}{cm} et \SI{1,0}{cm}.
\end{enumerate}

\section{Drapeaux américains}
La Lune est observée à l'aide de la plus grande lunette astronomique du monde, celle de Yerkes (États-Unis) qui possède un objectif de \SI{102}{cm} de diamètre et de \SI{19,2}{cm} de distance focale. Elle permet d'atteindre un grossissement maximal de \num{2500}.\\
Si l'angle sous lequel l'\oe il observe deux points est inférieur au pouvoir séparateur de l'\oe il $\epsilon = \SI{3E-4}{\radian }$, alors les deux points sont confondus.
\begin{enumerate}
    \item Déterminer la taille du plus petit détail que l'on peut voir sur la Lune en utilisant la lunette de Yerkes.
    \item En déduire si les drapeaux laissés par les Américains lors de leurs différents voyages sur la Lune sont visibles. (Taille d'un drapeau $\approx \SI{1}{m}$)
\end{enumerate}
\end{multicols}
\end{document}